% Хавсралт A
%\chapter{Хавсралтын нэр} % Main appendix title
\label{AppendixA} % For referencing this appendix elsewhere, use \ref{AppendixA}
\chapter*{Чухал кодын хэсгүүд}

\section*{Хэрэглэгчийн бүртгэл үүсгэх кодын хэсэг}
\paragraph{} Бүртгэл үүсгэх API endpoint /US-070/ нь Django REST Framework -ийн \texttt{generics.CreateAPIView} -ийг өргөтгөн хэрэгжүүлсэн. Хэрэглэгч имейл, нууц үгээ оруулсны дараа bcrypt -ээр hash хийн хадгалж, JWT \texttt{access} болон \texttt{refresh} токеныг буцаана.
\begin{lstlisting}[language=Python]
from rest_framework import generics, status, permissions
from rest_framework.response import Response
from rest_framework_simplejwt.tokens import RefreshToken
from .models import User
from .serializers import RegisterSerializer, UserProfileSerializer


class RegisterView(generics.CreateAPIView):
    """POST /api/auth/register/  -- US-070"""
    queryset           = User.objects.all()
    serializer_class   = RegisterSerializer
    permission_classes = [permissions.AllowAny]

    def create(self, request, *args, **kwargs):
        s = self.get_serializer(data=request.data)
        s.is_valid(raise_exception=True)
        user    = s.save()
        refresh = RefreshToken.for_user(user)
        return Response({
            "user":    UserProfileSerializer(user).data,
            "refresh": str(refresh),
            "access":  str(refresh.access_token),
        }, status=status.HTTP_201_CREATED)
\end{lstlisting}

\section*{Имэйл, нууц үгээр нэвтрэх кодын хэсэг}
\paragraph{} Нэвтрэх API endpoint нь \texttt{authenticate()} функцээр имэйл, нууц үгийг шалгаж, амжилттай тохиолдолд JWT токен буцаана. \texttt{is\_banned} статустай хэрэглэгчид 403 алдаа буцаана.
\begin{lstlisting}[language=Python]
from rest_framework.views import APIView
from django.contrib.auth import authenticate

class LoginView(APIView):
    """POST /api/auth/login/  -- US-070"""
    permission_classes = [permissions.AllowAny]

    def post(self, request):
        email    = request.data.get("email")
        password = request.data.get("password")
        user     = authenticate(request, username=email, password=password)
        if not user:
            return Response({"error": "Invalid credentials"},
                            status=status.HTTP_401_UNAUTHORIZED)
        if user.is_banned:
            return Response({"error": "Account is banned"},
                            status=status.HTTP_403_FORBIDDEN)
        refresh = RefreshToken.for_user(user)
        return Response({
            "user":    UserProfileSerializer(user).data,
            "refresh": str(refresh),
            "access":  str(refresh.access_token),
        })
\end{lstlisting}

\section*{Crowd Aggregation алгоритмын кодын хэсэг}
\paragraph{} Олон хэрэглэгчийн оруулсан сегментийн саналуудыг нэгтгэн \emph{давамгай} нөхцлийг тооцоолох алгоритм /US-030, US-031/. Тухайн сегментийн орчинд (radius $\approx 0.001$ градус) байрлах бүх саналыг тоолж green / yellow / red тус бүрийн саналаас хамгийн ихийг сонгоно.
\begin{lstlisting}[language=Python]
import hashlib
from django.db import models


class CrowdAggregation(models.Model):
    """Aggregated crowd vote per map cell -- E4 US-030 US-031"""
    DOMINANT_CHOICES = [
        ("green",  "Green"),
        ("yellow", "Yellow"),
        ("red",    "Red"),
        ("none",   "No data"),
    ]
    segment_hash = models.CharField(max_length=64, unique=True, db_index=True)
    green_votes  = models.PositiveIntegerField(default=0)
    yellow_votes = models.PositiveIntegerField(default=0)
    red_votes    = models.PositiveIntegerField(default=0)
    dominant     = models.CharField(max_length=10,
                                    choices=DOMINANT_CHOICES, default="none")

    @staticmethod
    def make_hash(start_lat, start_lng, end_lat, end_lng):
        key = (f"{round(float(start_lat),3)},{round(float(start_lng),3)},"
               f"{round(float(end_lat),3)},{round(float(end_lng),3)}")
        return hashlib.sha256(key.encode()).hexdigest()[:64]

    def compute_dominant(self):
        """
        US-031: most votes wins.
        Example: green=10, yellow=3, red=6 -> green
        """
        votes = {
            "green":  self.green_votes,
            "yellow": self.yellow_votes,
            "red":    self.red_votes,
        }
        total = sum(votes.values())
        if total == 0:
            self.dominant = "none"
        else:
            self.dominant = max(votes, key=votes.get)
        return self.dominant
\end{lstlisting}

\section*{Сегментийн санал нэгтгэх trigger функцийн кодын хэсэг}
\paragraph{} Сегмент шинээр хадгалагдах эсвэл шинэчлэгдэх бүрд автоматаар дуудагдах функц. Тухайн координатын ойролцоох (bounding box) бүх сегментийг тоолон \texttt{CrowdAggregation} хүснэгтийг шинэчилнэ.
\begin{lstlisting}[language=Python]
from .models import CrowdAggregation

def update_aggregation(segment):
    """
    Called after a Segment is saved.
    Re-counts all votes for the same spatial hash
    and recalculates dominant. -- E4 US-030
    """
    from apps.segments.models import Segment as Seg

    seg_hash = CrowdAggregation.make_hash(
        segment.start_lat, segment.start_lng,
        segment.end_lat,   segment.end_lng,
    )

    # Find all segments in the same spatial cell
    nearby = Seg.objects.filter(
        start_lat__range=(float(segment.start_lat) - 0.001,
                          float(segment.start_lat) + 0.001),
        start_lng__range=(float(segment.start_lng) - 0.001,
                          float(segment.start_lng) + 0.001),
        end_lat__range  =(float(segment.end_lat)   - 0.001,
                          float(segment.end_lat)   + 0.001),
        end_lng__range  =(float(segment.end_lng)   - 0.001,
                          float(segment.end_lng)   + 0.001),
    )

    g = nearby.filter(condition="green").count()
    y = nearby.filter(condition="yellow").count()
    r = nearby.filter(condition="red").count()

    agg, _ = CrowdAggregation.objects.get_or_create(
        segment_hash=seg_hash,
        defaults={
            "start_lat": segment.start_lat,
            "start_lng": segment.start_lng,
            "end_lat":   segment.end_lat,
            "end_lng":   segment.end_lng,
        }
    )
    agg.green_votes  = g
    agg.yellow_votes = y
    agg.red_votes    = r
    agg.compute_dominant()
    agg.save()
    return agg
\end{lstlisting}

\section*{POI байрлуулах болон санал өгөх кодын хэсэг}
\paragraph{} Хэрэглэгч газрын зураг дээр POI үүсгэх (US-020), upvote / downvote санал өгөх (US-022) ба админ POI -г батлах эсвэл татгалзах (US-051) логикийг агуулсан \texttt{ViewSet}.
\begin{lstlisting}[language=Python]
from rest_framework import viewsets, permissions
from rest_framework.decorators import action
from rest_framework.response import Response
from django.db import transaction
from .models import POI, POIVote
from .serializers import POISerializer


class POIViewSet(viewsets.ModelViewSet):
    """
    POST /api/pois/             -- create POI, status=pending  (US-020)
    POST /api/pois/{id}/vote/   -- upvote / downvote           (US-022)
    POST /api/pois/{id}/approve/-- moderator approval          (US-051)
    POST /api/pois/{id}/reject/ -- moderator rejection         (US-051)
    """
    serializer_class = POISerializer
    filterset_fields = ["poi_type", "status"]

    def perform_create(self, serializer):
        poi  = serializer.save(user=self.request.user, status="pending")
        user = self.request.user
        user.total_pois += 1
        user.save(update_fields=["total_pois"])

    @action(detail=True, methods=["post"])
    def vote(self, request, pk=None):
        """POST /api/pois/{id}/vote/  -- US-022 (one vote per user)"""
        if not request.user.is_authenticated:
            return Response({"error": "Authentication required"}, status=401)
        poi       = self.get_object()
        vote_type = request.data.get("vote_type")
        if vote_type not in ("up", "down"):
            return Response({"error": "vote_type must be up or down"}, status=400)

        with transaction.atomic():
            existing = POIVote.objects.filter(poi=poi, user=request.user).first()
            if existing:
                if existing.vote_type == vote_type:
                    # Toggle off
                    if vote_type == "up":
                        poi.upvotes = max(0, poi.upvotes - 1)
                    else:
                        poi.downvotes = max(0, poi.downvotes - 1)
                    existing.delete()
                else:
                    # Switch vote
                    if vote_type == "up":
                        poi.upvotes  += 1
                        poi.downvotes = max(0, poi.downvotes - 1)
                    else:
                        poi.downvotes += 1
                        poi.upvotes    = max(0, poi.upvotes - 1)
                    existing.vote_type = vote_type
                    existing.save()
            else:
                POIVote.objects.create(poi=poi, user=request.user,
                                       vote_type=vote_type)
                if vote_type == "up":
                    poi.upvotes += 1
                else:
                    poi.downvotes += 1
            poi.save()

        return Response({"upvotes": poi.upvotes,
                         "downvotes": poi.downvotes,
                         "user_vote": vote_type})

    @action(detail=True, methods=["post"])
    def approve(self, request, pk=None):
        """POST /api/pois/{id}/approve/  -- US-051"""
        poi = self.get_object()
        poi.status = "approved"
        poi.save()
        return Response({"status": "approved", "id": poi.id})

    @action(detail=True, methods=["post"])
    def reject(self, request, pk=None):
        """POST /api/pois/{id}/reject/  -- US-051"""
        poi    = self.get_object()
        reason = request.data.get("reason", "")
        if not reason:
            return Response({"error": "Rejection reason is required"},
                            status=400)
        poi.status        = "rejected"
        poi.reject_reason = reason
        poi.save()
        return Response({"status": "rejected"})
\end{lstlisting}
